\documentclass{zhslides}

\usepackage{logic}


\title{Day 4: 说谎者悖论}
%\subtitle{}
\date{\today}
\date{2026年7月23日}
\author{于佳铭}
\institute{唯理书院2026\\精确的逻辑与模糊的意义}

\begin{document}

\maketitle

%\begin{frame}{目录}
%  \tableofcontents
%\end{frame}

% 可以用 \alert{} 来高亮



\begin{frame}
  
  \begin{quote}
    这句话是假的。
  \end{quote}
  
\end{frame}



\section{说谎者悖论}

\begin{frame}{真理理论}
  
  塔斯基的形式化真理理论：
  
  \begin{law*}{T-schema}
    语句{``$P$\,\textquotedblright}为真当且仅当$P$。
  \end{law*}
  \vspace{4ex}
  \pause
  
  \begin{quote}
    $(L)$\quad $L$是假的。
  \end{quote}
  \pause
  应用T-schema得到：$L$为真，当且仅当$L$是假的。\contradiction
  
\end{frame}


\begin{frame}{塔斯基的语言层级}
  
  \bit
    \item 说谎者悖论来源于语言的自指性
    \item 放弃自然语言的封闭性，区分\term{元语言}{metalanguage}和\term{对象语言}{object language}
  \eit
  \vspace{4ex}
  \pause
  
  \begin{quote}
    'Schnee ist wei\ss' is true if and only if snow is white.
  \end{quote}
  
\end{frame}


\begin{frame}[t]{多值逻辑与说谎者悖论}

  \vspace{4ex}
  \begin{brainstorm*}{}
    三值逻辑能否解答说谎者悖论？
  \end{brainstorm*}
  \vspace{2ex}
  \pause
  
  \bit
    \item 克里普克(Kripke)：真值\term{缝隙}{gap}
    \item 说谎者句子在这个理论下不处于任何层级，因此真值不被确定——非真也非假
  \eit
  \vspace{2ex}
  \pause
  
  复仇悖论：
  \begin{quote}
    $(L^+)$\quad $L^+$不是真的。
  \end{quote}

\end{frame}



\section{双面真理论}

\begin{frame}{爆炸原理}
  
  \begin{law*}{矛盾律, LNC}
    对于任意语句 $\phi$，$\neg(\phi \land \neg \phi)$ 是永真式。
  \end{law*}
  \vspace{1ex}
  \pause
  
  一旦矛盾律不被满足，根据爆炸原理就可以得出一切结论为真。
  
  \begin{law*}{爆炸原理, Principle of explosion}
    对于任意命题$P$和$Q$，如果$P$和$\neg P$同时为真，那么可以推导出$Q$为真。
  \end{law*}
  这就是著名的拉丁短语``\textit{Ex falso quodlibet}'' (from falsehood, anything follows)
  
\end{frame}


\begin{frame}{双面真理论}
  
  \bit
    \item 普里斯特（Priest）：说谎者句子是真值\term{溢出}{glut}而非缝隙
    \item 放弃矛盾律，转而使用\term{次协调逻辑}{paraconsistent logic}，例如普利斯特的Logic of Paradox (LP)
  \eit
  
\end{frame}









\end{document}